\paragraph{\trLefschetzFibration{}。}\label{sub:lefschetz-fibrations}
设 $S$ 是一个\trConnected\trCompact、带非空边界的\trRiemann%
\trSurface。选取一个 $1$-形式 $\theta_S$，使得 $\theta_S|_{\partial S}>0$ 且
$\phi_S=d\theta_S>0$，以此把 $S$ 做成一个\trCorneredBoundary{}的%
\trExactSymplectic\trManifold。一个（\trExact{}的）\trLefschetzFibration，
就是一个\trFibrationWithSingularities{}$\pi\colon E^{2n+2}\to S$，其\trCriticalPoint{}是\trGeneric{}的
（也称\trNondegenerate）而且局部可积。这意味着：$I_E$ 在 $\mathrm{Crit}(\pi)$
的一个邻域内可积，并且 $D\pi$（视为由复\trLinear{}映射组成的丛
$\Hom_\CC(TE,\pi^*TS)$ 的一张\trSection）与\trZeroSection{}\trTransverse。
第二个条件与如下说法相\trEquivalence：复\trHessian{}$D^2\pi$ 在每个%
\trCriticalPoint{}处都是\trNondegenerate{}的。此外我们还将假设：每条\trFiber{}中至多只有一个%
\trCriticalPoint，从而\trProjection{}$\mathrm{Crit}(\pi)\to\mathrm{Critv}(\pi)$
是\trBijective{}的；这最后一条假设只是为了方便，其实很容易去掉，代价是让后
文的记号变得略显笨拙。
\trNondegenerate{}性蕴含诸\trCriticalPoint{}都是\trIsolated{}的，于是
$\mathrm{Crit}(\pi)$ 是 $\mathrm{int}(E)$ 的有限子集，同样地，$\mathrm{Critv}(\pi)$
是 $\mathrm{int}(S)$ 的有限子集。在每个\trCriticalPoint{}及其值附近，都有%
\trHolomorphic{}坐标，使 $\pi$ 在其中成为标准二次映射
\begin{equation}\label{eq:standard-quadratic-map}
Q(x)=x_1^2+\cdots+x_{n+1}^2
\end{equation}
一般来说，$\phi_E$ 在这些坐标里并不是标准的。然而，我们可以找到这一%
\trFibration{}的一个\trDeformation，它沿 $\partial^h E$ \trWellBehaved{}（而且实际
上在\trCriticalPoint{}附近是局部的），使得在\trDeformation{}的另一端，
\trKaehler{}形式在给定的\trHolomorphic\trMorse\trChart{}里成为标准形式。这就愈发增加了
下面这个基本模型的重要性：
\begin{example}\label{thm:basic-local-model}
设 $Q\colon\CC^{n+1}\to\CC$ 如 \eqref{eq:standard-quadratic-map} 中那样，而
$k\colon\CC^{n+1}\to\RR^+$ 是函数 $k(x)=(|x|^4-|Q(x)|^2)/4$。取定 $r,s>0$，
定义
\begin{equation}\label{eq:basic-model-domain}
E=\bigl\{x\in\CC^{n+1} : |Q(x)|\le r,\ k(x)\le s\bigr\},
\end{equation}
并给它配备 $\CC^{n+1}$ 上标准\trSymplecticForm{}的限制、其标准原函数
$(i/4)(z\,d\bar{z}-\bar{z}\,dz)$、给定的\trComplexStructure{}$I_E=i$，以及限
制 $Q$ 所得的映射 $\pi\colon E\to S=rD$。\trBoundaryFace{}为
\begin{equation*}
\partial^v E=\bigl\{x\in E : |\pi(x)|=r\bigr\},\qquad
\partial^h E=\bigl\{x\in E : k(x)=s\bigr\}.
\end{equation*}
选取截断函数 $k$ 的用意，是使 $TE^h$ 平行于 $\partial^h E$。为看出这一点，
注意 $TE^h_x$ 在 $\CC$ 上由 $(\nabla Q)_x=2\bar{x}$ 生成，并验证
$dk_x(\bar{x})=0$、$dk_x(i\bar{x})=0$。每一条\trNonsingular\trFiber{}$E_z$（$z\neq 0$）
都辛\trIsomorphic{}于 $B_s^*S^n\subset T^*S^n$ 中由长度（按标准\trMetric）$\le s$ 的%
\trCotangent{}向量组成的子集。明确地说，$B_s^*S^n=\{(u,v)\in\RR^{n+1}\times
S^n : \langle u,v\rangle=0\}$，配以\trSymplecticForm{}$du\wedge dv$；而对
$z>0$，\trIsomorphism{}$E_z\to B_s^*S^n$ 由下式给出：
\begin{equation*}
\phi_z(x)=\bigl(-\mathrm{im}(x)|\mathrm{re}(x)|,\,
\mathrm{re}(x)|\mathrm{re}(x)|^{-1}\bigr).
\end{equation*}
遗憾的是，虽然负\trLiouville{}（负\trRadial{}）\trVectorField{}确实沿
$\partial E$ 指向内侧，$E$ 却并不是弱 $I_E$-\trConvex{}的（诸\trFiber{}是，
\trTotalSpace{}不是）。因此，这还算不上这里所定义的\trExact\trLefschetzFibration{}的
一个例子——尽管从纯粹辛的观点看，它具备全部想要的性质。当然，可以修改
$I_E$ 来改善局面，但那样做并没有什么实际意义，因为归根结底 $E$ 只是一个局
部模型而已。
\end{example}
\begin{remark}\label{thm:other-lefschetz-definitions}
关于辛\trLefschetzFibration{}的其他定义，见例如 [39] 或 [64]。通常底与\trFiber{}
都被假定为\trClosed{}的，这在技术上稍稍简单一些。局部正规形式
\eqref{eq:standard-quadratic-map} 的存在性是\trHolomorphic\trMorseLemma{} [8] 的一
个推论（更精确地说，其陈述是：对底上任意选定的一组\trHolomorphic{}坐标，
都能找到\trTotalSpace{}上的坐标使 $\pi=Q$）。能使\trSymplecticStructure{}在这种坐标里成为标准的
那个\trDeformation，其构造见 [131, Lemma 1.6]。\cref{thm:basic-local-model}
的几何在 [131, Section 1.2] 中有穷尽的讨论。
\end{remark}
